\chapter{NLP Architectures}
\label{chap:nlp_architectures}

\begin{tldr}
This chapter covers the three major Transformer paradigms for NLP: encoder-only (BERT family), decoder-only (GPT family), and encoder-decoder (T5/BART). Understanding when to use each architecture—and the design decisions within each family—is a common staff-level interview topic. We also cover tokenization methods (BPE, WordPiece, SentencePiece), modern LLM innovations, and practical guidance for architecture selection.
\end{tldr}

\section{The Three Transformer Paradigms}
\label{sec:transformer_paradigms}

\subsection{Architectural Overview}

\begin{figure}[H]
\centering
\begin{tikzpicture}[
    node distance=0.8cm,
    box/.style={rectangle, draw, minimum width=2.5cm, minimum height=0.7cm, rounded corners},
    encoder/.style={box, fill=lightblue},
    decoder/.style={box, fill=lightgreen},
    cross/.style={box, fill=orange!30}
]
    % Encoder-only
    \begin{scope}[shift={(0,0)}]
        \node[encoder] (e1) {Encoder};
        \node[encoder, above=0.3cm of e1] (e2) {Encoder};
        \node[above=0.3cm of e2] (edots) {$\vdots$};
        \node[encoder, above=0.3cm of edots] (eN) {Encoder};
        \node[below=0.5cm of e1] {Input};
        \node[above=0.5cm of eN] {[CLS] / Tokens};
        \node[below=1.2cm of e1, font=\bfseries] {Encoder-Only};
        \node[below=1.7cm of e1, font=\small] {(BERT)};
        
        % Bidirectional arrows
        \draw[<->, thick, blue] ($(e1.west)+(-0.3,0)$) -- ($(eN.west)+(-0.3,0)$);
        \node[left=0.5cm of e2, font=\tiny, blue] {Bidirectional};
    \end{scope}
    
    % Decoder-only
    \begin{scope}[shift={(5,0)}]
        \node[decoder] (d1) {Decoder};
        \node[decoder, above=0.3cm of d1] (d2) {Decoder};
        \node[above=0.3cm of d2] (ddots) {$\vdots$};
        \node[decoder, above=0.3cm of ddots] (dN) {Decoder};
        \node[below=0.5cm of d1] {Input};
        \node[above=0.5cm of dN] {Next Token};
        \node[below=1.2cm of d1, font=\bfseries] {Decoder-Only};
        \node[below=1.7cm of d1, font=\small] {(GPT)};
        
        % Causal arrow
        \draw[->, thick, green!70!black] ($(d1.east)+(0.3,0)$) -- ($(dN.east)+(0.3,0)$);
        \node[right=0.5cm of d2, font=\tiny, green!70!black] {Causal};
    \end{scope}
    
    % Encoder-Decoder
    \begin{scope}[shift={(10,0)}]
        \node[encoder] (ed_e1) {Encoder};
        \node[encoder, above=0.3cm of ed_e1] (ed_eN) {Encoder};
        
        \node[decoder, right=1cm of ed_e1] (ed_d1) {Decoder};
        \node[cross, above=0.3cm of ed_d1] (ed_cross) {Cross-Attn};
        \node[decoder, above=0.3cm of ed_cross] (ed_dN) {Decoder};
        
        \node[below=0.5cm of ed_e1] {Source};
        \node[below=0.5cm of ed_d1] {Target};
        \node[above=0.5cm of ed_dN] {Output};
        \node[below=1.2cm of $(ed_e1)!0.5!(ed_d1)$, font=\bfseries] {Encoder-Decoder};
        \node[below=1.7cm of $(ed_e1)!0.5!(ed_d1)$, font=\small] {(T5, BART)};
        
        % Cross attention arrow
        \draw[->, thick, orange] (ed_eN.east) -- (ed_cross.west);
    \end{scope}
\end{tikzpicture}
\caption{The three Transformer paradigms for NLP.}
\label{fig:transformer_paradigms}
\end{figure}

\subsection{Key Differences}

\begin{table}[H]
\centering
\caption{Comparison of Transformer paradigms}
\begin{tabularx}{\textwidth}{lXXX}
\toprule
\textbf{Aspect} & \textbf{Encoder-Only} & \textbf{Decoder-Only} & \textbf{Encoder-Decoder} \\
\midrule
Attention & Bidirectional & Causal (left-to-right) & Bidirectional enc, Causal dec \\
Context & Full sequence & Past tokens only & Full source, past target \\
Pretraining & Masked LM (MLM) & Causal LM & Span corruption / denoising \\
Output & Representations & Next token prob & Next token prob \\
Best for & Understanding & Generation & Seq-to-seq \\
Examples & BERT, RoBERTa & GPT, LLaMA & T5, BART \\
\bottomrule
\end{tabularx}
\end{table}

%==============================================================================
\section{Encoder-Only Models (BERT Family)}
\label{sec:encoder_only}
%==============================================================================

\subsection{BERT: Bidirectional Encoder Representations from Transformers}

\subsubsection{Architecture}

\begin{itemize}
    \item Stack of Transformer encoder blocks
    \item Bidirectional self-attention (no causal mask)
    \item Special tokens: [CLS] for classification, [SEP] for segment separation
    \item Segment embeddings for sentence pairs
\end{itemize}

\subsubsection{Pre-training Objectives}

\textbf{Masked Language Modeling (MLM)}:
\begin{itemize}
    \item Randomly mask 15\% of tokens
    \item Of masked positions:
    \begin{itemize}
        \item 80\%: Replace with [MASK]
        \item 10\%: Replace with random token
        \item 10\%: Keep original
    \end{itemize}
    \item Predict original token
\end{itemize}

\begin{equation}
\mathcal{L}_{MLM} = -\E\left[\sum_{i \in \mathcal{M}} \log p(x_i | \tilde{\vx})\right]
\end{equation}

where $\mathcal{M}$ is the set of masked positions and $\tilde{\vx}$ is the corrupted input.

\textbf{Why the 80/10/10 split?}
\begin{itemize}
    \item [MASK] never appears at fine-tuning time
    \item Training only on [MASK] creates mismatch
    \item Random replacement forces model to maintain representations for all tokens
    \item Original token prevents trivial copying
\end{itemize}

\textbf{Next Sentence Prediction (NSP)}:
\begin{itemize}
    \item Given [CLS] Sentence A [SEP] Sentence B [SEP]
    \item Binary classification: Is B the actual next sentence?
    \item 50\% positive (actual next), 50\% negative (random)
\end{itemize}

Note: NSP was later found to be less important; RoBERTa removes it.

\subsubsection{Fine-tuning Patterns}

\begin{figure}[H]
\centering
\begin{tikzpicture}[
    node distance=0.5cm,
    box/.style={rectangle, draw, minimum width=1.5cm, minimum height=0.5cm, rounded corners, font=\small},
    bert/.style={box, fill=lightblue, minimum width=4cm}
]
    % Classification
    \begin{scope}[shift={(0,0)}]
        \node[bert] (bert1) {BERT};
        \node[box, fill=lightgreen, above=0.3cm of bert1] (cls1) {Classifier};
        \node[below=0.3cm of bert1, font=\small] {[CLS] sent [SEP]};
        \node[above=0.3cm of cls1, font=\small] {Label};
        \draw[->] (bert1) -- (cls1);
        \node[below=1cm of bert1, font=\bfseries\small] {Classification};
    \end{scope}
    
    % NER
    \begin{scope}[shift={(5,0)}]
        \node[bert] (bert2) {BERT};
        \node[box, fill=lightgreen, above=0.3cm of bert2, minimum width=4cm] (ner) {Token Classifier};
        \node[below=0.3cm of bert2, font=\small] {$t_1$ $t_2$ ... $t_n$};
        \node[above=0.3cm of ner, font=\small] {$l_1$ $l_2$ ... $l_n$};
        \draw[->] (bert2) -- (ner);
        \node[below=1cm of bert2, font=\bfseries\small] {Token Classification};
    \end{scope}
    
    % QA
    \begin{scope}[shift={(10,0)}]
        \node[bert] (bert3) {BERT};
        \node[box, fill=lightgreen, above=0.3cm of bert3, minimum width=4cm] (qa) {Start/End Classifier};
        \node[below=0.3cm of bert3, font=\small] {[CLS] Q [SEP] Context};
        \node[above=0.3cm of qa, font=\small] {Start, End positions};
        \draw[->] (bert3) -- (qa);
        \node[below=1cm of bert3, font=\bfseries\small] {Extractive QA};
    \end{scope}
\end{tikzpicture}
\caption{Common BERT fine-tuning patterns.}
\label{fig:bert_finetuning}
\end{figure}

\subsection{BERT Variants: What Mattered}

The 2019--2021 encoder race produced a long list of variants; for interview purposes, each survives as one transferable idea. The table below is the working summary---the per-model deep dives (RoBERTa's ablations, ALBERT's factorization math, DeBERTa's disentangled-attention equations) are canonical in the NLP volume [NLP~5] and are rarely worth their page count in a 2026 loop.

\begin{table}[H]
\centering
\caption{BERT family: the one idea each variant contributed}
\begin{tabularx}{\textwidth}{lrXX}
\toprule
\textbf{Model} & \textbf{Params} & \textbf{Key change} & \textbf{Lasting lesson} \\
\midrule
BERT (2018) & 110M / 340M & MLM + NSP pretraining & Bidirectional pretraining transfers \\
RoBERTa (2019) & 125M / 355M & Drop NSP; dynamic masking; 10$\times$ data, bigger batches & BERT was undertrained---recipe $>$ architecture \\
ALBERT (2019) & 12M--235M & Factorized embeddings; cross-layer weight sharing & Parameter count $\neq$ capacity; sharing trades quality for memory \\
ELECTRA (2020) & 14M--335M & Replaced-token detection instead of MLM & Objective design controls sample efficiency (see box) \\
DeBERTa (2020) & 100M--1.5B & Disentangled content/position attention; enhanced mask decoder & Relative-position modeling; long-running NLU SOTA (SuperGLUE) \\
ModernBERT (2024) & 139M / 395M & LLM-era recipe backported: RoPE, GeGLU, alternating local--global attention, 8K context, 2T tokens & Encoders benefit from the same modernization as decoders; the 2026 default for new encoder deployments \\
\bottomrule
\end{tabularx}
\end{table}

\begin{keyinsight}[ELECTRA: Signal Density Determines Sample Efficiency]
MLM wastes most of each batch: only the 15\% masked positions produce gradient signal. ELECTRA instead trains a small generator to fill in masks and a discriminator to classify \emph{every} token as original-or-replaced---so all positions contribute signal. Result: BERT-level quality at roughly $1/4$ the pretraining compute; the discriminator (not the generator) is kept for fine-tuning. The general principle---denser training signal per token wins---is the same argument that favors causal LM over MLM at scale, and it recurs whenever pretraining objectives are compared.
\end{keyinsight}

\subsection{When to Use Encoder-Only}

\begin{itemize}
    \item \textbf{Text classification}: Sentiment, topic, intent
    \item \textbf{Token classification}: NER, POS tagging
    \item \textbf{Extractive QA}: Find answer span in context
    \item \textbf{Semantic similarity}: Sentence embeddings, retrieval
    \item \textbf{NOT for}: Open-ended generation
\end{itemize}

%==============================================================================
\section{Decoder-Only Models (GPT Family)}
\label{sec:decoder_only}
%==============================================================================

\subsection{Architecture}

\begin{itemize}
    \item Stack of Transformer decoder blocks (no cross-attention)
    \item Causal self-attention (each position only attends to past)
    \item No [CLS] token; use last token or pooling for classification
    \item Autoregressive generation
\end{itemize}

\subsection{Causal Language Modeling}

\begin{mathresult}{Causal LM Objective}
\begin{equation}
\mathcal{L}_{CLM} = -\sum_{i=1}^{n} \log p(x_i | x_1, \ldots, x_{i-1}; \theta)
\end{equation}
\end{mathresult}

\textbf{Key property}: Every token provides a training signal (vs. 15\% for MLM).

\textbf{Inference}: Generate token by token, feeding each output back as input.

\subsection{Model Evolution}

\begin{table}[H]
\centering
\caption{GPT family evolution}
\begin{tabularx}{\textwidth}{lllX}
\toprule
\textbf{Model} & \textbf{Params} & \textbf{Context} & \textbf{Key Innovation} \\
\midrule
GPT (2018) & 117M & 512 & First large-scale CLM pretraining + task fine-tuning \\
GPT-2 (2019) & 1.5B & 1024 & Zero-shot task transfer from raw WebText (40GB) \\
GPT-3 (2020) & 175B & 2048 & In-context / few-shot learning at scale \\
GPT-4 (2023) & undisclosed & 8K/32K (later 128K) & Multimodal input; RLHF-aligned frontier quality \\
o-series (o1/o3, 2024--25) & undisclosed & 128K+ & RL-trained reasoning models: long chain-of-thought optimized with verifiable rewards, quality scaling with test-time compute (Chapter~\ref{chap:reinforcement_learning}) \\
\bottomrule
\end{tabularx}
\end{table}

After GPT-3, OpenAI stopped disclosing parameter counts, architectures, and training data; treat any specific figure you have heard for GPT-4-class models as rumor, not fact---quoting one in an interview signals the opposite of insider knowledge. The honest table therefore stops tracking sizes and tracks \emph{capability regimes} instead.

\subsection{Open-Weight LLMs}

The architecture column below is the one to internalize: the attention variant (MHA $\to$ GQA $\to$ MLA) determines KV-cache size, and dense-vs-MoE determines the total-vs-active parameter split---together they drive the serving economics discussed in Chapter~\ref{chap:inference_optimization}.

\begin{table}[H]
\centering
\small
\caption{Major open-weight LLMs (params shown as total/active for MoE models)}
\begin{tabularx}{\textwidth}{@{}>{\raggedright\arraybackslash}p{2.5cm}>{\raggedright\arraybackslash}p{2.1cm}>{\raggedright\arraybackslash}p{1.4cm}p{1.0cm}>{\raggedright\arraybackslash}p{2.2cm}>{\raggedright\arraybackslash}X@{}}
\toprule
\textbf{Model} & \textbf{Params} & \textbf{Ctx} & \textbf{Attn} & \textbf{FFN} & \textbf{Key Features} \\
\midrule
LLaMA (2023) & 7--65B & 2K & MHA & Dense & Modern recipe: RMSNorm, SwiGLU, RoPE \\
Llama 2 (2023) & 7--70B & 4K & GQA\textsuperscript{*} & Dense & GQA at 70B; open license \\
Mistral 7B (2023) & 7B & 8K & GQA & Dense & Sliding-window attn (v0.1 only) \\
Mixtral (2023) & 47B/13B & 32K & GQA & MoE, 8 experts, top-2 & First strong open MoE \\
Llama 3.1 (2024) & 8--405B & 128K & GQA & Dense & 15T tokens; 128K vocab \\
Gemma 2 (2024) & 2--27B & 8K & GQA & Dense & Interleaved local--global attention \\
DeepSeek-V2 (2024) & 236B/21B & 128K & MLA & MoE fine-grained + shared & MLA shrinks KV cache $\sim$10$\times$ \\
Qwen2.5 (2024) & 0.5--72B & 128K & GQA & Dense & Full size ladder; strong multilingual + code \\
DeepSeek-V3 (2024) & 671B/37B & 128K & MLA & MoE, 256 experts, top-8 + shared & Aux-loss-free balancing; MTP; FP8 training \\
Gemma 3 (2025) & 1--27B & 128K & GQA & Dense & 5:1 local:global; multimodal \\
Qwen3 (2025) & 0.6--32B; 30B/3B; 235B/22B & 32--128K & GQA & Dense + MoE & Hybrid thinking / non-thinking modes \\
Llama 4 (2025) & 109B/17B; 400B/17B & 10M (claimed) & GQA & MoE & Natively multimodal; no-RoPE interleaved layers \\
\bottomrule
\end{tabularx}

{\raggedright\footnotesize \textsuperscript{*}GQA in the 34B/70B sizes; the 7B/13B sizes remain MHA.\par}
\end{table}

\subsection{Modern LLM Design Choices}

\subsubsection{LLaMA Architecture Details}

\begin{enumerate}
    \item \textbf{Pre-normalization}: RMSNorm before attention and FFN
    \item \textbf{SwiGLU activation}: $\text{FFN}(x) = (\text{Swish}(xW_1) \odot xV)W_2$
    \item \textbf{RoPE}: Rotary positional embeddings
    \item \textbf{No bias}: Remove bias terms throughout
    \item \textbf{GQA} (LLaMA 2): Grouped-query attention for efficiency
\end{enumerate}

\subsubsection{Mistral Innovations (v0.1)}

\begin{enumerate}
    \item \textbf{Sliding window attention (SWA)}: Each layer attends to a window of 4096 tokens
    \item \textbf{Rolling buffer cache}: Fixed memory regardless of sequence length
    \item \textbf{Pre-fill and chunking}: Efficient handling of long prompts
\end{enumerate}

Note: Mistral dropped SWA after v0.1---later versions use full attention over longer contexts. SWA survives as one point in the sparse-attention design space, notably in interleaved local--global hybrid form (e.g., Gemma 2/3 alternate sliding-window and full-attention layers).

\subsection{Mixture of Experts in the Decoder Stack}
\label{subsec:moe_decoder}

Most 2026 frontier models---DeepSeek-V3/R1, Llama 4, Qwen3's large sizes, Mixtral, and (by broad consensus) the closed frontier models---are \textbf{Mixture of Experts (MoE)} decoders. The change is localized: attention layers are untouched, but each FFN is replaced by $N$ expert FFNs plus a small learned \textit{router}. For every token, the router (a linear layer over the hidden state) scores all experts and dispatches the token to the top-$k$; the layer's output is the router-weighted sum of those $k$ expert outputs. Classic designs used few large experts ($N{=}8$, $k{=}2$ in Mixtral); the fine-grained style uses many small ones plus always-on shared experts ($N{=}256$, $k{=}8$, one shared in DeepSeek-V3).

This splits the parameter count into two numbers you must keep separate:
\begin{itemize}
    \item \textbf{Total parameters} --- everything, all experts included. Determines quality ceiling and \emph{memory footprint}: every expert must be resident to serve, because you cannot predict which experts the next token needs. DeepSeek-V3: 671B total.
    \item \textbf{Active parameters} --- the subset touched by one token (attention + router + $k$ experts + shared). Determines \emph{per-token FLOPs}, and therefore decode latency and serving cost. DeepSeek-V3: 37B active.
\end{itemize}

\begin{keyinsight}[Why the Frontier Went MoE]
At fixed training compute, a sparse model reaches lower loss than a dense one: capacity scales with total parameters while cost scales with active parameters. At serving time the same asymmetry pays again---a 671B/37B model decodes with roughly the per-token compute of a 37B dense model. The price is paid in memory and systems complexity: the weights no longer fit one GPU (expert parallelism with all-to-all communication), and small-batch serving leaves most experts idle. MoE trades cheap capacity (HBM) for expensive compute (FLOPs)---exactly the right trade at datacenter scale, and usually the wrong one on a single consumer GPU.
\end{keyinsight}

For interviews, be ready for the standard probe: ``Mixtral is 8$\times$7B---why is it 47B total, not 56B, and what does 13B active mean for serving?'' (Experts replace only FFNs, so attention/embeddings are shared, giving 46.7B; top-2 routing touches $\sim$12.9B per token, so latency resembles a 13B model while memory requires holding all 46.7B.) Everything deeper---router training dynamics, auxiliary load-balancing losses and their failure modes, expert-choice routing, DeepSeek's fine-grained + shared expert design and V3's aux-loss-free balancing, and expert-parallel serving---is treated canonically in Chapter~\ref{chap:efficient_architectures}; this section exists so the decoder-stack story is complete without it.

\subsection{In-Context Learning}

\textbf{In-Context Learning.} The ability to perform tasks by conditioning on examples in the prompt, without gradient updates.

\textbf{Example (few-shot sentiment)}:
\begin{verbatim}
Review: "Great movie!" Sentiment: Positive
Review: "Terrible waste of time" Sentiment: Negative
Review: "Absolutely loved it!" Sentiment: [MODEL GENERATES]
\end{verbatim}

\textbf{Emergent at scale}: ICL doesn't appear in small models; emerges around 10B+ parameters.

\subsection{Generation Strategies}

\subsubsection{Greedy Decoding}
\begin{equation}
x_t = \argmax_x p(x | x_{<t})
\end{equation}
Simple but repetitive, lacks diversity.

\subsubsection{Beam Search}
Maintain $k$ best partial sequences.
\begin{itemize}
    \item Better for constrained generation (translation)
    \item Still tends toward repetition for open-ended generation
\end{itemize}

\subsubsection{Temperature Sampling}
\begin{equation}
p(x | x_{<t}) = \frac{\exp(z_x / T)}{\sum_{x'} \exp(z_{x'} / T)}
\end{equation}

\begin{itemize}
    \item $T < 1$: Sharper, more deterministic
    \item $T > 1$: Flatter, more random
    \item $T = 1$: Standard softmax
\end{itemize}

\subsubsection{Top-k Sampling}
Sample only from top-$k$ most likely tokens.

\subsubsection{Nucleus (Top-p) Sampling}
Sample from smallest set of tokens whose cumulative probability $\geq p$.

\begin{equation}
V_p = \argmin_V \left\{\sum_{x \in V} p(x | x_{<t}) \geq p\right\}
\end{equation}

\textbf{Advantage}: Adapts to probability distribution shape (more tokens when uncertain).

\subsubsection{Min-p Sampling}
Keep only tokens whose probability is at least a fraction $p_{\min}$ of the \emph{top} token's probability:

\begin{equation}
V_{\min\text{-}p} = \left\{x : p(x | x_{<t}) \geq p_{\min} \cdot \max_{x'} p(x' | x_{<t})\right\}
\end{equation}

The truncation threshold scales with the model's confidence: when the top token has probability 0.9, only near-certain alternatives survive; when the distribution is flat, many candidates pass. This fixes top-p's failure mode at high temperature---once temperature flattens the distribution, a fixed cumulative mass $p$ admits a long tail of junk tokens, whereas min-p stays anchored to the best token. Typical $p_{\min} = 0.05$--$0.1$; widely adopted in open-source serving stacks (llama.cpp, vLLM) since 2024, and the practical default for creative generation at $T > 1$.

\subsubsection{Recommended Settings}

\begin{table}[H]
\centering
\caption{Generation settings by use case}
\begin{tabularx}{\textwidth}{lXX}
\toprule
\textbf{Use Case} & \textbf{Settings} & \textbf{Rationale} \\
\midrule
Factual QA & $T=0$ or greedy & Deterministic, accurate \\
Creative writing & $T=0.7-1.0$, top-p=0.9 & Diverse but coherent \\
Code generation & $T=0.2-0.5$, top-p=0.95 & Low diversity, high precision \\
Brainstorming & $T=1.0+$, top-p=0.95 & Maximum diversity \\
High-temperature creative & $T=1.0-1.5$, min-p=0.05-0.1 & Diversity without tail junk \\
\bottomrule
\end{tabularx}
\end{table}

\subsubsection{Beyond Token-by-Token Sampling}

\textbf{Speculative decoding.} Autoregressive decode is memory-bandwidth-bound: each token requires reading all the weights to produce one token. Speculative decoding exploits this by letting a cheap \textit{draft} model propose several tokens, which the target model then verifies in a single parallel forward pass; a rejection-sampling correction accepts a prefix of the draft such that the output distribution is \emph{provably identical} to sampling from the target model alone. Accepted drafts amortize the weight reads over multiple tokens, giving a typical 2--3$\times$ decode speedup with zero quality change---the closest thing to a free lunch in serving. Draft-model selection, self-drafting variants (Medusa/EAGLE-style extra heads), and the interaction with batching are covered in Chapter~\ref{chap:inference_optimization}.

\textbf{Structured (constrained) decoding.} Production systems increasingly require output that parses: JSON matching a schema, function-call arguments, SQL. Constrained decoding compiles the target grammar or JSON schema into a finite-state machine over the \emph{tokenizer vocabulary}; at each step, tokens that cannot lead to a valid continuation are masked out of the logits before sampling, guaranteeing syntactically valid output with negligible overhead. Two caveats: token/character boundary alignment is subtle (a grammar over characters must be lifted to a vocabulary of multi-character tokens), and hard masks can force the model off its preferred phrasing---overly rigid schemas measurably degrade content quality. Implementations (Outlines/XGrammar-style engines and their integration into vLLM/TGI) are likewise covered in Chapter~\ref{chap:inference_optimization}.

\subsection{When to Use Decoder-Only}

\begin{itemize}
    \item \textbf{Text generation}: Stories, articles, completions
    \item \textbf{Conversational AI}: Chatbots, assistants
    \item \textbf{Code generation}: Copilot-style applications
    \item \textbf{Few-shot learning}: When you can't fine-tune
    \item \textbf{General-purpose}: When one model must do many tasks
\end{itemize}

%==============================================================================
\section{The Anatomy of Parameters}
\label{sec:parameter_anatomy}
%==============================================================================

Model tables quote sizes; interviews test whether you can \emph{derive} them. For a standard decoder block with model dimension $d$, the parameter budget decomposes cleanly, and knowing the decomposition lets you reconstruct any model's size from its config---or spot a wrong config in seconds.

\subsection{The Per-Layer Budget}

\textbf{Attention: $4d^2$.} Four projection matrices $\mW^Q, \mW^K, \mW^V, \mW^O$, each $d \times d$ under multi-head attention (the $h$ heads partition $d$; they do not multiply the parameter count). With GQA, K and V shrink to $d \times (n_{kv} \cdot d_{head})$, reducing attention to $2d^2 + 2 d \cdot n_{kv} d_{head}$---about $2.25$--$2.5d^2$ in typical 8-KV-head configs. A useful second-order fact: that saves only $\sim$10\% of the \emph{block's} parameters (GQA's real win is KV-cache size, not weights), so $4d^2$ remains a fine mental estimate.

\textbf{FFN: $\approx 8d^2$.} The classic two-matrix FFN with hidden size $4d$ costs $d \times 4d + 4d \times d = 8d^2$. SwiGLU uses \emph{three} matrices ($\mW_{gate}$ and $\mW_{up}$, each $d \times d_{ff}$, plus $\mW_{down}$, $d_{ff} \times d$), so designers shrink the hidden size to $d_{ff} \approx \tfrac{2}{3} \cdot 4d = \tfrac{8}{3}d$ to hold the budget at $\approx 8d^2$. This is why Llama configs show the otherwise-mysterious $d_{ff} = 11008 \approx 2.7d$ rather than $4d$.

\textbf{Everything else: negligible.} RMSNorm adds $2d$ per layer; modern models drop all biases. Rounding error.

\begin{mathresult}{Transformer Parameter Count}
\begin{equation}
P \;\approx\; \underbrace{12\, L\, d^2}_{\text{blocks: } 4d^2 \text{ attn} + 8d^2 \text{ FFN}} \;+\; \underbrace{V d \times \{1 \text{ tied},\; 2 \text{ untied}\}}_{\text{embeddings + LM head}}
\end{equation}
where $L$ is the number of layers, $d$ the model dimension, and $V$ the vocabulary size.
\end{mathresult}

\subsection{Worked Sanity Check: Llama-2-7B}

Config: $L = 32$, $d = 4096$, $d_{ff} = 11008$, $V = 32000$, untied embeddings.

\begin{table}[H]
\centering
\caption{Where Llama-2-7B's parameters live}
\begin{tabular}{llrr}
\toprule
\textbf{Component} & \textbf{Formula} & \textbf{Count} & \textbf{Share} \\
\midrule
Attention & $4d^2 \times L = 4 \cdot 4096^2 \cdot 32$ & 2.15B & 32\% \\
FFN (SwiGLU) & $3 d\, d_{ff} \times L = 3 \cdot 4096 \cdot 11008 \cdot 32$ & 4.33B & 64\% \\
Embeddings + head & $2 V d = 2 \cdot 32000 \cdot 4096$ & 0.26B & 4\% \\
\midrule
Total & & \textbf{6.74B} & 100\% \\
\bottomrule
\end{tabular}
\end{table}

The quick version---$12 \cdot 32 \cdot 4096^2 = 6.4$B plus $0.26$B of embeddings $\approx 6.7$B---lands within 1\% of the true 6.74B. Two portable takeaways: roughly \textbf{two-thirds of a dense transformer is FFN} (which is exactly why MoE targets the FFN, Section~\ref{subsec:moe_decoder}), and embedding share depends on scale and vocabulary: negligible at 7B with a 32K vocab, but $\sim$1.05B ($13\%$) for Llama-3-8B's 128K untied vocab, and roughly a third of GPT-2-small. Small multilingual models spend a startling fraction of themselves on the embedding table.

This section covers the weight side only: the activation-side arithmetic (KV-cache bytes per token and what GQA/MLA do to it) lives in Chapter~\ref{chap:inference_optimization}, and the full transformer derivation is canonical in the NLP volume [NLP~5].

\begin{interviewq}
\textbf{Q: You're told a model is ``Llama-class, 7B parameters.'' Roughly where do those 7B parameters live?}

\badgeLfive{L5} \quad \badgetype{Estimation} \quad \badgefreq{\freqhigh}

\stronganswer Start from the per-layer budget. A decoder block is attention plus FFN: attention is four $d \times d$ projections $= 4d^2$; the FFN is $\approx 8d^2$ (classic $4d$ hidden, or SwiGLU's three matrices at $d_{ff} \approx \tfrac{8}{3}d$). So $\approx 12d^2$ per layer. Llama-7B has $d = 4096$ and $L = 32$: $12 \cdot 32 \cdot 4096^2 \approx 6.4$B. Embeddings add $2Vd = 2 \cdot 32000 \cdot 4096 \approx 0.26$B (untied input embedding + LM head), giving $\approx 6.7$B---the actual model is 6.74B.

The split: $\sim$64\% FFN, $\sim$32\% attention, $\sim$4\% embeddings. Worth stating unprompted: the FFN dominance is why MoE replaces FFNs specifically, and why FFN-heavy quantization/pruning schemes exist.

\textit{Adjustments an interviewer may probe:} GQA trims only the attention slice ($4d^2 \to \sim 2.5d^2$; roughly 10\% of the block---its real win is KV cache, not parameters). A 128K vocabulary pushes embeddings to $\sim$1B---material at 8B scale, dominant below 1B. Tied embeddings halve the embedding term. For MoE, the formula splits: total params replicate the FFN term per expert, active params count only the routed $k$.

\redflags
\begin{itemize}
    \item Assuming attention dominates---the FFN is $2\times$ larger
    \item Forgetting the LM head / not knowing whether embeddings are tied
    \item Unable to explain why $d_{ff} = 11008$ rather than $4d = 16384$ (SwiGLU's three-matrix budget)
    \item Needing the config handed over---the point of the question is reconstructing it
\end{itemize}

\interviewtest{Whether your familiarity with LLMs is load-bearing: can you regenerate the numbers in the model card from first principles, or have you only read them?}

\goodgreat
\begin{itemize}
    \item \badgeLfive{L5} Produces $12Ld^2$ + embeddings and the right FFN/attention/embedding split
    \item \badgeLsix{L6} Handles the SwiGLU and GQA corrections, does the 6.7B sanity check unprompted, and connects FFN share to MoE design
    \item \badgeLseven{L7} Extends fluently to total-vs-active MoE accounting, embedding-share scaling across model sizes and vocabularies, and the weights-vs-KV-cache memory split at serving time
\end{itemize}

\followups{``How does the split change for Mixtral 8$\times$7B---derive the 47B.'' ``Why is $d_{ff} \approx 2.7d$ in Llama?'' ``At what model size does a 256K vocabulary become the dominant cost?''}
\end{interviewq}

%==============================================================================
\section{Encoder-Decoder Models (T5 Family)}
\label{sec:encoder_decoder}
%==============================================================================

\subsection{Architecture}

\begin{itemize}
    \item Separate encoder and decoder stacks
    \item Encoder: Bidirectional self-attention
    \item Decoder: Causal self-attention + cross-attention to encoder
    \item Both encoder and decoder outputs used
\end{itemize}

\subsection{T5: Text-to-Text Transfer Transformer}

\subsubsection{Key Innovation: Unified Text-to-Text Format}

All tasks framed as text-to-text:
\begin{itemize}
    \item \textbf{Translation}: ``translate English to German: [text]'' → German text
    \item \textbf{Classification}: ``sentiment: [text]'' → ``positive'' / ``negative''
    \item \textbf{Summarization}: ``summarize: [text]'' → summary
    \item \textbf{QA}: ``question: [q] context: [c]'' → answer
\end{itemize}

\subsubsection{Pre-training: Span Corruption}

\begin{enumerate}
    \item Randomly select spans (average 3 tokens each)
    \item Replace each span with sentinel token ([MASK1], [MASK2], ...)
    \item Decoder outputs: sentinel + original span tokens
\end{enumerate}

\begin{example}
Input: ``The [MASK1] brown fox [MASK2] the lazy dog''\\
Target: ``[MASK1] quick [MASK2] jumped over''
\end{example}

\subsubsection{T5 Variants}

\begin{table}[H]
\centering
\caption{T5 model sizes}
\begin{tabular}{lrrr}
\toprule
\textbf{Model} & \textbf{Params} & \textbf{Layers} & \textbf{Dim} \\
\midrule
T5-Small & 60M & 6 enc + 6 dec & 512 \\
T5-Base & 220M & 12 + 12 & 768 \\
T5-Large & 770M & 24 + 24 & 1024 \\
T5-3B & 3B & 24 + 24 & 1024 \\
T5-11B & 11B & 24 + 24 & 1024 \\
\bottomrule
\end{tabular}
\end{table}

\subsection{BART: Denoising Autoencoder}

\subsubsection{Pre-training: Multiple Corruption Schemes}

\begin{enumerate}
    \item \textbf{Token masking}: Replace with [MASK]
    \item \textbf{Token deletion}: Remove tokens (no placeholder)
    \item \textbf{Text infilling}: Replace span with single [MASK]
    \item \textbf{Sentence permutation}: Shuffle sentences
    \item \textbf{Document rotation}: Rotate to start from random token
\end{enumerate}

\textbf{Best configuration}: Text infilling with 30\% span masking, Poisson span lengths ($\lambda = 3$).

\subsection{When to Use Encoder-Decoder}

\begin{itemize}
    \item \textbf{Translation}: Explicit alignment between source and target
    \item \textbf{Summarization}: Condense long source to short target
    \item \textbf{Structured generation}: When output structure differs from input
    \item \textbf{Conditional generation}: Strong conditioning on input needed
\end{itemize}

%==============================================================================
\section{Architecture Selection Guide}
\label{sec:nlp_arch_selection}
%==============================================================================

\begin{figure}[H]
\centering
\begin{tikzpicture}[
    node distance=1cm,
    decision/.style={diamond, draw, aspect=2.5, inner sep=1pt, fill=lightyellow, font=\small},
    result/.style={rectangle, draw, rounded corners, fill=lightgreen, minimum width=2cm, font=\small},
    arrow/.style={->, thick}
]
    % Start
    \node[decision] (gen) {Need generation?};
    
    % No generation
    \node[decision, below left=1.5cm and 1cm of gen] (und) {Task type?};
    \node[result, below left=1cm and 0cm of und] (bert) {BERT/RoBERTa};
    \node[result, below right=1cm and 0cm of und] (deberta) {DeBERTa};
    
    % Generation
    \node[decision, below right=1.5cm and 1cm of gen] (cond) {Conditional?};
    \node[result, below left=1cm and 0cm of cond] (t5) {T5/BART};
    \node[result, below right=1cm and 0cm of cond] (gpt) {GPT/LLaMA};
    
    % Edges
    \draw[arrow] (gen) -- node[above left, font=\small] {No} (und);
    \draw[arrow] (gen) -- node[above right, font=\small] {Yes} (cond);
    \draw[arrow] (und) -- node[left, font=\small] {Classification} (bert);
    \draw[arrow] (und) -- node[right, font=\small] {SOTA needed} (deberta);
    \draw[arrow] (cond) -- node[left, font=\small] {Yes} (t5);
    \draw[arrow] (cond) -- node[right, font=\small] {No/Few-shot} (gpt);
\end{tikzpicture}
\caption{NLP architecture selection decision tree.}
\label{fig:nlp_decision_tree}
\end{figure}

\begin{table}[H]
\centering
\caption{Architecture selection by task}
\begin{tabularx}{\textwidth}{lXX}
\toprule
\textbf{Task} & \textbf{Recommended} & \textbf{Alternative} \\
\midrule
Sentiment classification & BERT, RoBERTa & Fine-tuned LLM \\
Named Entity Recognition & BERT + CRF & Token classifier \\
Extractive QA & BERT, DeBERTa & - \\
Generative QA & T5, GPT & - \\
Machine translation & T5, mBART & - \\
Summarization & BART, T5 & GPT (zero-shot) \\
Dialogue / Chat & GPT, LLaMA & - \\
Code generation & CodeLLaMA, StarCoder & GPT-4 \\
Semantic similarity & Sentence-BERT & E5, BGE \\
Zero-shot classification & GPT, LLaMA & BART (NLI-based) \\
\bottomrule
\end{tabularx}
\end{table}

\begin{interviewq}
\textbf{Q: Encoder-only vs decoder-only vs encoder-decoder---when do you use each, and what drives the decision?}

\badgeLfive{L5} \quad \badgetype{Trade-off} \quad \badgefreq{\freqhigh}

\stronganswer The choice maps to three factors: task type, data availability, and latency requirements.

\textit{Encoder-only (BERT, RoBERTa, DeBERTa):} Use when the task produces fixed-size outputs from the full input context. Classification, NER, extractive QA, semantic similarity. Bidirectional attention means every token sees the full context, giving richer representations for understanding. Inference is a single forward pass---5--10$\times$ faster than generation. Best when: you have labeled data for fine-tuning, need low latency, and the task is ``understanding'' not ``generation.'' A fine-tuned DeBERTa-base is still the most cost-effective solution for text classification.

\textit{Decoder-only (GPT, LLaMA, Mistral):} Use when the task requires generating variable-length text, or when you want zero/few-shot capability. The ``prompt $\to$ completion'' interface is universal: classification becomes ``Classify: [text] Answer:'', translation becomes ``Translate: [text] Translation:''. Critical advantage: in-context learning eliminates the need for task-specific training data. Best when: you need generality across many tasks, have limited labeled data, or the output is open-ended. Downside: autoregressive generation is slow (sequential token-by-token) and expensive for simple tasks.

\textit{Encoder-decoder (T5, BART, mBART):} Use when the input and output have different structures and the input needs deep bidirectional understanding. Translation is the canonical example: the encoder deeply understands the source, cross-attention aligns source and target, and the decoder generates the target. Also strong for abstractive summarization and structured data-to-text. Best when: the task is genuinely sequence-to-sequence and you want the encoder's bidirectional context for the input. Downside: more parameters for the same ``effective'' model size, and the paradigm has been largely superseded by decoder-only models at scale.

\textit{Modern nuance:} At large scale (70B+), decoder-only models match or beat encoder-only models even on understanding tasks, because the massive parameter count compensates for the causal attention constraint. But at small scale (100M--1B), architecture choice still matters significantly.

\redflags
\begin{itemize}
    \item ``Always use GPT because it's the best''---ignores cost and latency
    \item Not knowing that encoder models are faster for classification
    \item Confusing when encoder-decoder adds value vs.\ when decoder-only suffices
\end{itemize}

\interviewtest{Practical architecture selection based on constraints, not just theoretical knowledge.}

\goodgreat
\begin{itemize}
    \item \badgeLfive{L5} Correctly maps task types to architecture families
    \item \badgeLsix{L6} Discusses the scale-dependent nature of the tradeoff, mentions cost and latency
    \item \badgeLseven{L7} Explains why decoder-only dominates at scale (training signal density, interface universality, simpler infrastructure), discusses when encoder models are still the right choice in production
\end{itemize}

\followups{``Would you use GPT-4 for production classification?'' ``Why did T5's encoder-decoder approach not scale as well as decoder-only?'' ``Can you use a decoder for embedding/retrieval tasks?''}
\end{interviewq}

%==============================================================================
\section{Tokenization}
\label{sec:tokenization}
%==============================================================================

\subsection{Why Tokenization Matters}

\begin{itemize}
    \item \textbf{Vocabulary size}: Affects embedding table size and rare word handling
    \item \textbf{Sequence length}: Character-level is longer, word-level shorter
    \item \textbf{OOV handling}: Subword methods handle unseen words gracefully
    \item \textbf{Multilingual}: Different scripts need different treatment
\end{itemize}

\subsection{Tokenization Methods}

\subsubsection{Byte Pair Encoding (BPE)}

\begin{algorithm}[H]
\caption{BPE Training}
\begin{algorithmic}[1]
\REQUIRE Corpus, vocabulary size $V$
\STATE Initialize vocabulary with all characters
\WHILE{$|$vocabulary$| < V$}
    \STATE Count frequency of all adjacent pairs
    \STATE Merge most frequent pair into new token
    \STATE Add new token to vocabulary
\ENDWHILE
\RETURN vocabulary, merge rules
\end{algorithmic}
\end{algorithm}

\textbf{Example}:
\begin{enumerate}
    \item Start: l, o, w, e, r, \_
    \item Most frequent pair: (e, r) → Merge to ``er''
    \item Next: (l, o) → ``lo''
    \item Continue until vocabulary size reached
\end{enumerate}

\textbf{Used in}: GPT-2, GPT-3, RoBERTa

\subsubsection{WordPiece}

Similar to BPE but uses likelihood-based merging:
\begin{equation}
\text{score}(x, y) = \frac{\text{freq}(xy)}{\text{freq}(x) \times \text{freq}(y)}
\end{equation}

Merge pair that maximizes likelihood of training data.

\textbf{Used in}: BERT, mBERT, DistilBERT

\textbf{Notation}: Continuations marked with \#\# (e.g., ``playing'' → ``play'', ``\#\#ing'')

\subsubsection{Unigram}

\begin{enumerate}
    \item Start with large vocabulary
    \item Iteratively remove tokens that least affect likelihood
    \item Keep tokens that maximize: $\sum_{\text{corpus}} \log p(\text{tokenization})$
\end{enumerate}

\textbf{Advantage}: Can output multiple tokenizations with probabilities.

\textbf{Used in}: T5, ALBERT

\subsubsection{SentencePiece}

\begin{itemize}
    \item Language-agnostic (no pre-tokenization)
    \item Treats input as raw byte stream
    \item Uses BPE or Unigram internally
    \item Handles any script (CJK, emoji, etc.)
\end{itemize}

\textbf{Used in}: LLaMA-1/2, T5, mT5

\subsection{Comparison}

\begin{table}[H]
\centering
\caption{Tokenization method comparison}
\begin{tabularx}{\textwidth}{lXXX}
\toprule
\textbf{Method} & \textbf{Pros} & \textbf{Cons} & \textbf{Used By} \\
\midrule
BPE & Efficient, deterministic & Greedy & GPT, RoBERTa \\
WordPiece & Likelihood-based & More complex & BERT \\
Unigram & Probabilistic, multiple tokenizations & Slower training & T5, ALBERT \\
SentencePiece & Language-agnostic, handles any script & — & LLaMA, T5 \\
\bottomrule
\end{tabularx}
\end{table}

\subsection{Practical Considerations}

\begin{itemize}
    \item \textbf{Vocabulary size}: 30K--50K in the BERT/GPT-2 era; modern LLMs run 32K--256K (Llama 2: 32K; Llama 3: 128K; Gemma: 256K)---larger vocabularies buy shorter sequences and better multilingual fertility at the cost of embedding parameters (Section~\ref{sec:parameter_anatomy})
    \item \textbf{Byte-level BPE}: GPT-2 uses byte-level (256 base tokens) for complete coverage
    \item \textbf{Special tokens}: [PAD], [UNK], [CLS], [SEP], [MASK] in the BERT era; modern chat models instead reserve control tokens for conversation structure (Section~\ref{subsec:chat_templates}), e.g.\ \verb+<|im_start|>+ and \verb+<|eot_id|>+
    \item \textbf{Normalization}: Lowercasing, unicode normalization, etc.
\end{itemize}

%==============================================================================
\section{Instruction Tuning and RLHF}
\label{sec:instruction_tuning}
%==============================================================================

\subsection{The Alignment Problem}

Base language models are trained to predict next tokens, not to be helpful. They may:
\begin{itemize}
    \item Continue harmful text
    \item Refuse benign requests
    \item Not follow instructions
    \item Hallucinate confidently
\end{itemize}

\subsection{Instruction Tuning}

Fine-tune on instruction-following examples:
\begin{verbatim}
Instruction: Summarize the following article in one sentence.
Input: [article text]
Output: [one-sentence summary]
\end{verbatim}

\textbf{Effect}: Model learns to follow diverse instructions without task-specific fine-tuning.

\textbf{Data, then and now.} The 2022--23 generation---FLAN and Natural Instructions (academic NLP tasks converted to instructions), Alpaca and Dolly (tens of thousands of self-instruct or crowdsourced examples)---established the idea but is no longer what competitive post-training looks like. The 2025-era open reference points are: \textit{Tulu-3-style recipes}---fully documented multi-stage pipelines (curated + synthetic SFT at the $\sim$1M-example scale, then preference optimization, then RL on verifiable rewards) with decontaminated prompt mixtures and per-stage evals; \textit{UltraFeedback-class preference data}---tens of thousands of prompts with multiple model completions scored by a strong LLM judge, which replaced expensive human preference labels for most open work; and \textit{synthetic SFT at millions scale}---persona-conditioned generation and reasoning traces distilled from frontier models, quality-filtered rather than hand-written. The through-line: data quality, diversity, and decontamination now matter more than raw instruction count, and ``instruction tuning'' is just the first stage of a longer post-training pipeline---covered end-to-end (RLHF, DPO variants, GRPO, RLVR) in Chapter~\ref{chap:reinforcement_learning}.

\subsection{Chat Templates and Special Tokens}
\label{subsec:chat_templates}

A base checkpoint models raw text. An instruct/chat model is the same architecture trained on conversations rendered through a \textbf{chat template}: a fixed serialization in which reserved control tokens delimit roles and turns. The template is part of the model artifact---as load-bearing as the weights. ChatML-style templates (GPT/Qwen lineage) look like:

\begin{verbatim}
<|im_start|>system
You are a helpful assistant.<|im_end|>
<|im_start|>user
What is GQA?<|im_end|>
<|im_start|>assistant
\end{verbatim}

Llama 3 uses the same idea with different furniture: \verb+<|begin_of_text|>+, role headers like \verb+<|start_header_id|>+\allowbreak\verb+user<|end_header_id|>+, and \verb+<|eot_id|>+ ending each turn. Three practical facts:

\begin{itemize}
    \item \textbf{Control tokens are single vocabulary items}, typically reserved at tokenizer-training time but only given \emph{meaningful embeddings during post-training}. To a base model they are near-random vectors; to the instruct model they are the strongest structural signal in the sequence.
    \item \textbf{The system prompt has a template-defined location.} The model was trained to treat text in that slot as standing instructions; a system prompt pasted into the user turn (or omitted where training always included one) measurably weakens instruction-following.
    \item \textbf{Turn-end $\neq$ end-of-text.} Chat models learn to emit a turn-end token (\verb+<|im_end|>+, \verb+<|eot_id|>+) when done; the serving stack must treat exactly that token as the stop condition. The document-level EOS is a different token with a different job.
\end{itemize}

\textbf{Why template mismatch silently destroys quality.} Fine-tune with serialization A and serve with serialization B (or none), and nothing crashes---the model just sees token sequences from outside its training distribution. Training loss even decreases normally, because the model happily learns your malformed format; the damage only shows at inference against the original template, or in production quality. Classic symptoms: ignored system prompts, role confusion, degraded instruction-following, and runaway generations that never emit the stop token. Classic causes: hand-rolled prompt strings instead of the tokenizer's bundled template (\texttt{apply\_chat\_template}), doubled BOS tokens (template adds one, framework adds another), loss masks that exclude the turn-end token so the model never learns to stop, and tokenizer version drift between training and serving.

\begin{warningbox}
Tokenizer and template mismatches between training and inference are silent killers: fine-tune with one tokenizer or template and serve with another (or another \emph{version}), and the model degrades or produces garbage with no error raised. Version-pin the tokenizer \emph{and} the chat template with the checkpoint, and CI-check them: serialize a fixed conversation through the serving stack and diff the token IDs against the training pipeline's output.
\end{warningbox}

Tokenization algorithms themselves are summarized in Section~\ref{sec:tokenization}, with the deep treatment canonical in the NLP volume [NLP~12].

\begin{interviewq}
\textbf{Q: Your fine-tuned chat model performs worse than the base model it was tuned from. Walk me through your diagnosis.}

\badgeLsix{L6} \quad \badgetype{Debugging} \quad \badgefreq{\freqhigh}

\stronganswer Fine-tuning making a model \emph{worse} is almost always a serialization or bookkeeping bug, not a modeling problem---so I check the artifact before the algorithm. My ordered differential:

\textit{1. Chat-template mismatch (most common).} Training serialized conversations one way; serving serializes them another way (or not at all). Check: take one training example, run it through the \emph{serving} stack, and diff the exact token IDs against what the training pipeline produced. Look for hand-rolled prompt strings, doubled BOS (template inserts one, framework prepends another), and missing role headers. This 15-minute round-trip test catches most cases.

\textit{2. Wrong EOS / stop-token handling.} Symptom: generations that ramble and never stop, or truncate mid-sentence. Causes: the loss mask excluded the turn-end token, so the model never learned to emit it; or serving stops on the document EOS while the model emits \texttt{<|eot\_id|>}-style turn ends (or vice versa). Check: does the model ever produce the configured stop token? Is the turn-end token's loss unmasked in training?

\textit{3. Tokenizer drift.} Tokenizer library version changed between training and serving; special tokens were added and shifted IDs; the embedding matrix was resized during training but the old tokenizer files were shipped. Symptom ranges from subtle degradation to outright garbage. Check: assert identical vocab size, special-token map, and tokenizer version at load time.

\textit{4. Learning-rate-induced forgetting.} Only after 1--3 are excluded is this a training problem: too-high LR or too many epochs on narrow data overwrote general capability. Check: compare base vs.\ fine-tuned on general benchmarks (MMLU-style + a held-out instruction eval); a broad uniform drop implicates forgetting. Fix: lower LR ($\sim$1e-5 full fine-tune, or LoRA), fewer epochs, mix in general instruction data.

The unifying point: the deployable artifact is \emph{weights + tokenizer + template + stop-token config}, versioned together. Most ``fine-tuning failures'' are those four falling out of sync.

\redflags
\begin{itemize}
    \item Reaching for ``needs more data / bigger model'' before checking serialization
    \item Not knowing that training loss looks perfectly healthy under a wrong template---the model learns the malformed format just fine
    \item Treating EOS as one token with one job (document end vs.\ turn end)
    \item No concrete verification step (the token-ID diff)
\end{itemize}

\interviewtest{Production debugging reflexes for the single most common real-world fine-tuning failure class, and whether you understand a chat model as a versioned artifact rather than just weights.}

\goodgreat
\begin{itemize}
    \item \badgeLfive{L5} Knows template mismatch exists and suggests using \texttt{apply\_chat\_template}
    \item \badgeLsix{L6} Gives the ordered differential with concrete checks: token-ID round-trip diff, doubled BOS, loss-masked turn-end token, tokenizer version asserts
    \item \badgeLseven{L7} Frames the institutional fix---template/tokenizer as CI-checked, version-pinned artifacts; canary evals comparing base vs.\ tuned on general capability; knows why the failure is silent (loss decreases regardless)
\end{itemize}

\followups{``Why does training loss go down even with the wrong template?'' ``How would you detect a doubled BOS token in production?'' ``Does LoRA change the forgetting diagnosis?''}
\end{interviewq}

\subsection{RLHF: Reinforcement Learning from Human Feedback}

\begin{figure}[H]
\centering
\begin{tikzpicture}[
    node distance=1cm,
    box/.style={rectangle, draw, minimum width=2.5cm, minimum height=0.8cm, rounded corners},
    stage/.style={box, fill=lightblue}
]
    % Stage 1
    \node[stage] (sft) {SFT Model};
    \node[below=0.3cm of sft, font=\small] {Supervised Fine-Tuning};
    
    % Stage 2
    \node[stage, right=2cm of sft] (rm) {Reward Model};
    \node[below=0.3cm of rm, font=\small] {Train on preferences};
    
    % Stage 3
    \node[stage, right=2cm of rm] (ppo) {RLHF Model};
    \node[below=0.3cm of ppo, font=\small] {PPO optimization};
    
    % Edges
    \draw[->] (sft) -- (rm) node[midway, above, font=\small] {Generate};
    \draw[->] (rm) -- (ppo) node[midway, above, font=\small] {Score};
\end{tikzpicture}
\caption{RLHF training pipeline.}
\label{fig:rlhf_pipeline_overview}
\end{figure}

\subsubsection{Stage 1: Supervised Fine-Tuning (SFT)}

Fine-tune base model on high-quality demonstrations:
\begin{equation}
\mathcal{L}_{SFT} = -\sum_i \log p(y_i | x_i; \theta)
\end{equation}

\subsubsection{Stage 2: Reward Model Training}

Train model to predict human preferences:
\begin{equation}
\mathcal{L}_{RM} = -\E_{(x, y_w, y_l) \sim \mathcal{D}}\left[\log \sigma(r(x, y_w) - r(x, y_l))\right]
\end{equation}

where $y_w$ is preferred over $y_l$.

\subsubsection{Stage 3: RL Fine-Tuning}

The policy is trained to \emph{maximize} the KL-regularized reward objective:
\begin{equation}
J(\theta) = \E\left[r(x, y) - \beta \cdot \text{KL}(\pi_\theta(y|x) \| \pi_{ref}(y|x))\right]
\end{equation}

The KL penalty prevents the policy from diverging too far from the SFT model. Note that this is the objective, not the PPO loss itself: in practice, PPO---a clipped surrogate loss on importance ratios, paired with a learned value model for advantage estimation---is the mechanism used to ascend $J(\theta)$.

\subsection{Direct Preference Optimization (DPO)}

\textbf{Key insight}: Can optimize directly on preferences without explicit reward model.

\begin{equation}
\mathcal{L}_{DPO} = -\E\left[\log \sigma\left(\beta \log \frac{\pi_\theta(y_w|x)}{\pi_{ref}(y_w|x)} - \beta \log \frac{\pi_\theta(y_l|x)}{\pi_{ref}(y_l|x)}\right)\right]
\end{equation}

\textbf{Advantages over RLHF}:
\begin{itemize}
    \item No reward model needed
    \item No RL training instability
    \item Simpler implementation
    \item Often similar results
\end{itemize}

DPO (2023) is where this chapter's alignment story hands off: the 2024--26 continuation---online/iterative preference optimization, GRPO, RL with verifiable rewards, and reasoning-model training---is covered in Chapter~\ref{chap:reinforcement_learning}.

\begin{warningbox}
Don't assume bigger models are always better for your task. A fine-tuned BERT-base (110M params) often outperforms zero-shot GPT-4 on domain-specific classification tasks with even modest labeled data (1K+ examples). The right architecture depends on your data budget, latency requirements, and task type.
\end{warningbox}

\begin{interviewq}
\textbf{Q: Your fine-tuned LLM generates fluent but factually wrong answers. Diagnose the issue and propose solutions.}

\badgeLsix{L6} \quad \badgetype{Debugging} \quad \badgefreq{\freqhigh}

\stronganswer Fluent but factually incorrect output is the hallmark of ``hallucination''---the model generates plausible-sounding text that isn't grounded in reality. I would diagnose along three axes.

\textit{Diagnosis 1: Training data issues.} (a) The fine-tuning data may contain factual errors that the model memorized. Check: sample and manually verify fine-tuning examples. (b) The model may have been fine-tuned with too high a learning rate, degrading the factual knowledge stored in pre-trained weights (catastrophic forgetting). Check: compare base model vs.\ fine-tuned model on factual QA benchmarks. (c) The fine-tuning data may reward fluency over accuracy (e.g., RLHF reward model doesn't penalize factual errors).

\textit{Diagnosis 2: Decoding strategy.} High temperature or top-p sampling can push the model toward low-probability but fluent completions. For factual tasks, use lower temperature ($T \leq 0.3$) or greedy decoding. Check: does accuracy improve with $T=0$? If yes, the knowledge is there but sampling introduces errors.

\textit{Diagnosis 3: Knowledge boundary.} The model may be answering questions outside its training distribution. LLMs don't know what they don't know---they generate confidently regardless. Check: does the model hallucinate more on recent events, niche domains, or numerical questions?

\textit{Solutions:} (1) \textbf{RAG (Retrieval-Augmented Generation)}: Retrieve relevant documents and include them in context---grounds the model in actual data. This is the highest-impact intervention. (2) \textbf{Calibration training}: Fine-tune the model to say ``I don't know'' when uncertain, using examples where the answer isn't in its training data. (3) \textbf{Citation requirement}: Train the model to cite sources, which makes hallucinations detectable. (4) \textbf{Self-consistency}: Generate multiple answers and check agreement---inconsistency signals hallucination. (5) \textbf{Constitutional AI / RLHF with factuality reward}: Add factuality as an explicit reward signal.

\redflags
\begin{itemize}
    \item ``Just train on more data''---more data doesn't fix hallucination without targeted intervention
    \item Ignoring the decoding strategy as a contributing factor
    \item Not mentioning RAG as the primary practical solution
\end{itemize}

\interviewtest{Systematic debugging of LLM behavior, understanding of hallucination causes and mitigations.}

\goodgreat
\begin{itemize}
    \item \badgeLfive{L5} Identifies hallucination as the issue, suggests RAG
    \item \badgeLsix{L6} Provides structured diagnosis (data, decoding, knowledge boundary), multiple mitigations
    \item \badgeLseven{L7} Discusses the fundamental tension between fluency and factuality objectives, mentions process reward models for step-by-step verification, addresses evaluation of factuality at scale
\end{itemize}

\followups{``How do you measure hallucination rate?'' ``When does RAG fail?'' ``How does RLHF contribute to hallucination?''}
\end{interviewq}

\begin{interviewq}
\textbf{Q: Estimate the inference latency for generating 100 tokens with a 7B parameter model on a single A100 GPU.}

\badgeLsix{L6} \quad \badgetype{Estimation} \quad \badgefreq{\freqmed}

\stronganswer LLM inference has two phases: prefill (process the prompt) and decode (generate tokens one at a time). I'll estimate both.

\textit{Model memory.} 7B parameters in BF16 = $7 \times 10^9 \times 2$ bytes = 14 GB. Fits comfortably on a single A100 (80GB), leaving room for KV cache.

\textit{Prefill phase.} Assume a 500-token prompt. Prefill is compute-bound because we process all tokens in parallel. FLOPs $\approx 2 \times P \times n = 2 \times 7\text{B} \times 500 = 7 \times 10^{12}$ FLOPs. A100 BF16 throughput: 312 TFLOPS. Prefill time $\approx 7 \times 10^{12} / 312 \times 10^{12} \approx 22$ ms. In practice $\sim$30--50 ms accounting for overhead.

\textit{Decode phase.} Each token generation requires reading the full model weights from memory (memory-bandwidth bound, not compute-bound). A100 memory bandwidth: 2 TB/s. Model size: 14 GB. Time per token $\geq 14 \text{GB} / 2000 \text{GB/s} = 7$ ms. This is the theoretical minimum; in practice, accounting for KV cache reads and overhead, expect $\sim$10--15 ms per token. For 100 tokens: $100 \times 12$ ms $\approx$ 1.2 seconds.

\textit{Total latency:} Prefill ($\sim$40 ms) + decode ($\sim$1.2 s) $\approx$ \textbf{1.2--1.5 seconds} for 100 generated tokens. This gives $\sim$70--80 tokens/second for a single request.

\textit{Sanity check:} Benchmarks show LLaMA-2-7B on A100 achieves 50--80 tokens/s with optimized serving (vLLM, TGI), consistent with this estimate.

\redflags
\begin{itemize}
    \item Treating decode as compute-bound---it's memory-bandwidth bound for single-request serving
    \item Forgetting to distinguish prefill from decode phases
    \item Off by more than 5$\times$ on the final estimate
\end{itemize}

\interviewtest{Understanding of LLM inference bottlenecks and ability to do back-of-envelope calculations.}

\goodgreat
\begin{itemize}
    \item \badgeLfive{L5} Gets the right order of magnitude (seconds, not minutes or milliseconds)
    \item \badgeLsix{L6} Distinguishes prefill (compute-bound) from decode (memory-bound), uses memory bandwidth for per-token estimate
    \item \badgeLseven{L7} Discusses how batching changes the arithmetic (decode becomes compute-bound at batch size $\sim$64+), mentions quantization (INT4 halves per-token latency), discusses KV cache memory as the context-length bottleneck
\end{itemize}

\followups{``How does INT4 quantization change the numbers?'' ``What batch size makes decode compute-bound?'' ``How much memory does the KV cache need for 4K context?''}
\end{interviewq}

\begin{interviewq}
\textbf{Q: How does in-context learning work mechanistically? Why can large language models learn from examples in the prompt without any gradient updates?}

\badgeLseven{L7} \quad \badgetype{First Principles} \quad \badgefreq{\freqlow}

\stronganswer In-context learning (ICL) is one of the most surprising emergent properties of large LLMs. The model performs a new task given only a few examples in the prompt---no parameter updates. The leading mechanistic explanations are:

\textit{Theory 1: Implicit Bayesian inference.} The model has learned a distribution over tasks during pre-training. Given examples in the prompt, attention layers implement approximate Bayesian inference to identify the most likely task and apply the corresponding input-output mapping already encoded in the weights. This explains why example format matters as much as content---it helps the model ``locate'' the task.

\textit{Theory 2: Attention as mesa-optimization.} Akyurek et al.\ (2023) and Von Oswald et al.\ (2023) showed that transformer attention layers can implement a single step of gradient descent on the in-context examples. The key-value pairs act as ``training data'' and the attention operation computes something equivalent to a linear regression update. Multi-layer transformers implement multi-step optimization. This is ``learning to learn'' at a mechanistic level.

\textit{Theory 3: Task retrieval, not task learning.} The model doesn't learn new tasks from the prompt---it retrieves a pre-existing capability. The demonstrations serve as a ``task descriptor'' that activates the right circuit. Evidence: in earlier/smaller models, ICL performance was surprisingly robust to randomized demonstration labels (Min et al., 2022)---format seemed to matter more than content---and performance correlates with how well the task was represented in pre-training data. However, this result needs nuance: Wei et al.\ (2023) showed that larger models \emph{are} sensitive to label correctness and can even learn flipped label mappings from the prompt, overriding their semantic priors. So pure task retrieval cannot be the whole story at scale---large models genuinely use the input--label mappings in context.

\textit{When ICL breaks:} Tasks very different from pre-training distribution, tasks requiring novel reasoning not seen during training, and tasks where the label space is unfamiliar. ICL also has a strong recency bias (later examples in the prompt influence the output more) and is sensitive to example ordering.

\textit{Scale dependence:} ICL is nearly absent in models below $\sim$1B parameters and becomes reliable around 10B+. This suggests that the capacity to implement mesa-optimization requires sufficient model expressivity.

\redflags
\begin{itemize}
    \item ``The model fine-tunes on the examples''---there are no gradient updates
    \item Giving only one explanation without acknowledging the debate
    \item Not mentioning scale dependence
\end{itemize}

\interviewtest{Depth of understanding of emergent LLM capabilities and awareness of current research.}

\goodgreat
\begin{itemize}
    \item \badgeLfive{L5} Knows ICL exists and that it's about pattern matching, not gradient updates
    \item \badgeLsix{L6} Can articulate 2--3 theories, discusses practical implications (format sensitivity, ordering)
    \item \badgeLseven{L7} Knows the ``attention as gradient descent'' formalism, discusses scale dependence, and connects to the mesa-optimization framework
\end{itemize}

\followups{``Why does example order matter?'' ``Can ICL solve tasks not seen during pre-training?'' ``How do you decide between ICL and fine-tuning for a production task?''}
\end{interviewq}

\begin{interviewq}
\textbf{Q: BPE vs SentencePiece vs character-level tokenization---what are the tradeoffs and how does tokenizer choice affect model performance?}

\badgeLfive{L5} \quad \badgetype{Trade-off} \quad \badgefreq{\freqmed}

\stronganswer Tokenization is the first step in the NLP pipeline, and it profoundly affects sequence length, vocabulary size, rare word handling, and multilingual capability.

\textit{BPE (Byte Pair Encoding):} Starts with characters, iteratively merges the most frequent adjacent pair. Deterministic tokenization. Vocabulary: 30K--50K. Strengths: good balance between compression (short sequences) and coverage (handles unseen words via subwords). Weaknesses: greedy algorithm, no probabilistic model, and purely frequency-based merges may not capture linguistic structure. Used by GPT-2/3.

\textit{Byte-level BPE:} Variant that starts from 256 byte values instead of characters. Guarantees complete coverage of any input (no UNK tokens ever) without needing an explicit UNK token. Used by GPT-2/3 and Llama 3 (tiktoken-style vocabulary), and most current LLMs. Note that LLaMA-1/2 is \emph{not} byte-level BPE---it uses SentencePiece BPE with byte \emph{fallback}, where bytes are used only for characters missing from the vocabulary. Slight downside of byte-level: some non-English scripts get inefficient tokenization (one character may become 2--4 tokens).

\textit{SentencePiece:} Language-agnostic framework that treats input as raw bytes (no pre-tokenization). Implements both BPE and Unigram internally. Handles any script (CJK, Arabic, emoji) without language-specific preprocessing. The Unigram model starts with a large vocabulary and prunes tokens that least affect likelihood---this gives a probabilistic model that can output multiple tokenizations with probabilities. Used by T5 (Unigram mode) and LLaMA-1/2 (BPE mode with byte fallback).

\textit{Character-level:} Each character is a token. Vocabulary is tiny (256 for bytes), no OOV issues, but sequences become 3--5$\times$ longer, dramatically increasing compute cost (attention is $O(n^2)$). Works well for character-level tasks (spelling, morphology) but impractical for long-document processing. Used in some specialized models (ByT5).

\textit{Key tradeoff:} Vocabulary size vs.\ sequence length. Larger vocabulary $\to$ shorter sequences (faster) but larger embedding table (more parameters). Smaller vocabulary $\to$ longer sequences (slower) but better handling of rare/unseen words. The sweet spot for most LLMs is 32K--128K tokens with byte-level BPE.

\redflags
\begin{itemize}
    \item Treating tokenization as unimportant---it directly affects model capacity and efficiency
    \item Not knowing that BPE is greedy and SentencePiece can be probabilistic
    \item Ignoring multilingual considerations
\end{itemize}

\interviewtest{Understanding of tokenization fundamentals and their downstream effects on model design.}

\goodgreat
\begin{itemize}
    \item \badgeLfive{L5} Describes each method correctly, knows the vocab-vs-length tradeoff
    \item \badgeLsix{L6} Discusses byte-level BPE, multilingual efficiency, and the impact on arithmetic/reasoning tasks
    \item \badgeLseven{L7} Connects tokenization to ``fertility'' (tokens per word varies by language, creating unfairness), discusses tokenizer training data biases, mentions byte-level models (MegaByte) as a potential solution
\end{itemize}

\followups{``Why do LLMs struggle with character-level tasks like counting letters?'' ``How does vocabulary size affect model parameters?'' ``What happens if you fine-tune with a different tokenizer than pre-training?''}
\end{interviewq}

\begin{interviewq}
\textbf{Q: Why do decoder-only models dominate modern NLP despite encoder-decoder being theoretically more flexible?}

\badgeLsix{L6} \quad \badgetype{Conceptual} \quad \badgefreq{\freqhigh}

\stronganswer This is a crucial architectural question. Encoder-decoder models have a structural advantage for conditioned generation---the encoder can bidirectionally understand the input while the decoder generates the output. Yet decoder-only models (GPT, LLaMA, Claude) dominate. The reasons are primarily about scaling efficiency and interface design, not modeling capability.

\textit{Training efficiency.} In an encoder-decoder model, the encoder parameters only receive gradients from the decoder's cross-attention. The decoder must split its capacity between self-attention and cross-attention. A decoder-only model uses all its parameters for a single, unified objective. Be careful not to overstate this, though: controlled comparisons from the T5/UL2 era show encoder-decoder models are compute-competitive at matched FLOPs on many tasks, so ``decoder-only uses compute more efficiently'' is not a settled fact. The robust advantages are the ones that follow---interface, data format, and infrastructure.

\textit{Interface universality.} The ``text in, text out'' interface of decoder-only models handles every task through prompting: classification, translation, summarization, reasoning, code generation, and conversation. Encoder-decoder models work best for structured input-output tasks but are awkward for open-ended generation or conversation. One decoder-only model replaces dozens of specialized encoder-decoder models.

\textit{Simpler infrastructure.} Decoder-only models have one forward path: tokens in, next token out. Encoder-decoder requires managing two separate forward passes, cross-attention cache, and asymmetric batching. This complexity compounds across training, serving, and optimization. KV cache management is simpler for decoder-only (one cache vs.\ encoder output + decoder KV cache).

\textit{Scaling data.} Decoder-only models can be pre-trained on raw text with no special formatting. Encoder-decoder models need structured pairs (input $\to$ output) or corruption schemes (T5's span corruption). Raw text is vastly more abundant.

\textit{Where encoder-decoder still wins:} Machine translation (bidirectional understanding of source language is valuable), speech recognition (Whisper is encoder-decoder), and tasks with very long inputs and short outputs (summarization of books).

\redflags
\begin{itemize}
    \item ``Decoder-only is just objectively better''---the dominance is about scaling and infrastructure, not modeling
    \item Not acknowledging any advantage of encoder-decoder
    \item Confusing ``flexible'' with ``efficient''
\end{itemize}

\interviewtest{Understanding of the systems-level and scaling-level reasons behind architectural trends, beyond just modeling capabilities.}

\goodgreat
\begin{itemize}
    \item \badgeLfive{L5} Names the interface universality advantage
    \item \badgeLsix{L6} Discusses parameter efficiency, training data requirements, and infrastructure simplicity
    \item \badgeLseven{L7} Connects to the ``bitter lesson'' (simpler + scale > clever architecture), discusses prefix LM (decoder with bidirectional prefix) as a middle ground, mentions that Google's Gemini-family uses decoder-only despite having developed T5
\end{itemize}

\followups{``What is a prefix language model and does it get the best of both?'' ``Why is Whisper encoder-decoder while LLMs are decoder-only?'' ``Could decoder-only models ever be as good at translation as encoder-decoder?''}
\end{interviewq}

%==============================================================================
\section{Interview Questions}
\label{sec:nlp_interview}
%==============================================================================

\begin{interviewq}
\textbf{Q: Design the architecture for a multilingual document understanding system that handles 50+ languages, mixed-language documents, and both structured (tables, forms) and unstructured text.}

\badgeLseven{L7} \quad \badgetype{Architecture Design} \quad \badgefreq{\freqmed}

\stronganswer This is a complex system requiring multiple architectural decisions. I'll break it down by component.

\textit{Document preprocessing.} OCR pipeline (Tesseract or cloud OCR for text extraction from scans), layout analysis to identify text blocks, tables, headers, and figures. Use a layout-aware model like LayoutLMv3 or DocFormer that combines text tokens with spatial position information (bounding box coordinates) and optionally image patches.

\textit{Backbone model.} For multilingual capability, start with XLM-RoBERTa-large (550M params, trained on 100 languages) or mDeBERTa-v3. For document understanding specifically, LayoutLMv3 adds spatial embeddings: each token gets $(x_1, y_1, x_2, y_2)$ coordinates from OCR. This allows the model to understand that text in the same table column is related even if not adjacent in reading order.

\textit{Architecture design.}
\begin{enumerate}
    \item \textbf{Text + layout encoder}: LayoutLMv3 or similar, producing token-level representations with spatial awareness
    \item \textbf{Table-specific module}: For structured extraction, add a row/column attention mechanism that attends within table cells. Alternatively, serialize tables with special tokens (row/column separators)
    \item \textbf{Task heads}: Classification head for document type, extraction head for key-value pairs, NER head for entity extraction, span head for QA
\end{enumerate}

\textit{Multilingual handling.} XLM-RoBERTa handles code-switching (mixed languages) naturally because it was trained on multilingual data without language boundaries. For low-resource languages: leverage cross-lingual transfer---fine-tune on English (data-rich) and evaluate on target language. SentencePiece tokenizer with 250K vocabulary handles all scripts.

\textit{Scaling considerations.} For 50+ languages, data imbalance is the biggest challenge. Use temperature-based sampling during training: $p_L \propto n_L^{1/T}$ where $n_L$ is the number of examples in language $L$ and $T > 1$ upsamples low-resource languages. This prevents the model from overfitting to high-resource languages.

\textit{Serving.} Document processing is latency-tolerant (seconds, not milliseconds). Batch processing with model parallelism for large documents. Cache language detection to route to language-specific post-processing.

\redflags
\begin{itemize}
    \item Proposing separate models per language---doesn't scale to 50+
    \item Ignoring layout information for document understanding
    \item Not addressing low-resource languages
\end{itemize}

\interviewtest{System design thinking at the intersection of NLP, multilingual, and document AI.}

\goodgreat
\begin{itemize}
    \item \badgeLfive{L5} Chooses a multilingual model and mentions layout information
    \item \badgeLsix{L6} Designs the full pipeline with OCR, layout-aware model, and task-specific heads; addresses data imbalance
    \item \badgeLseven{L7} Discusses cross-lingual transfer, temperature sampling for language balance, serving architecture, and evaluation strategy across languages (not just English accuracy)
\end{itemize}

\followups{``How do you evaluate on languages you don't speak?'' ``What if a new language is added with zero training data?'' ``How does the system handle documents that mix tables and prose?''}
\end{interviewq}

\begin{interviewq}
\textbf{Q: Walk me through the DPO loss function. How does it eliminate the need for a separate reward model?}

\badgeLsix{L6} \quad \badgetype{First Principles} \quad \badgefreq{\freqhigh}

\stronganswer DPO (Direct Preference Optimization) is an alignment technique that trains directly on preference pairs without needing to train a reward model or run RL. Starting from the RLHF objective (maximize reward while staying close to reference policy), the key insight is that the optimal policy has a closed-form relationship to the reward function.

In RLHF, the optimal policy satisfies: $\pi^*(y|x) = \frac{1}{Z(x)} \pi_{\text{ref}}(y|x) \exp\left(\frac{r(x,y)}{\beta}\right)$. Rearranging to solve for $r$: $r(x,y) = \beta \log \frac{\pi^*(y|x)}{\pi_{\text{ref}}(y|x)} + \beta \log Z(x)$.

Substituting this into the Bradley-Terry preference model $p(y_w \succ y_l | x) = \sigma(r(x,y_w) - r(x,y_l))$, the $Z(x)$ terms cancel:

$$\loss_{\text{DPO}} = -\E_{(x, y_w, y_l)}\left[\log \sigma\left(\beta \log \frac{\pi_\theta(y_w|x)}{\pi_{\text{ref}}(y_w|x)} - \beta \log \frac{\pi_\theta(y_l|x)}{\pi_{\text{ref}}(y_l|x)}\right)\right]$$

\textit{Intuition:} DPO increases the log-probability ratio (vs.\ reference) for preferred outputs and decreases it for dispreferred outputs. The $\beta$ parameter controls how far the policy can drift from the reference. The implicit reward is the log-probability ratio itself.

\textit{Advantages over RLHF:} (1) No reward model to train (saves a training stage). (2) No RL instability (PPO is notoriously hard to tune). (3) Simpler implementation (just supervised fine-tuning with a special loss). (4) Often similar or better results in practice.

\textit{Limitations:} (1) Requires good preference data (garbage in, garbage out). (2) The reference model must be reasonably good (DPO refines, doesn't teach from scratch). (3) Can overfit to the preference dataset distribution. (4) Doesn't easily support online data collection (preference pairs are fixed).

\redflags
\begin{itemize}
    \item Not knowing the derivation from RLHF objective
    \item Saying DPO doesn't use a reward model at all---it has an \textit{implicit} reward model
    \item Claiming DPO always beats RLHF
\end{itemize}

\interviewtest{Mathematical understanding of alignment techniques and ability to derive results from first principles.}

\goodgreat
\begin{itemize}
    \item \badgeLfive{L5} Knows DPO simplifies RLHF by eliminating the reward model, can state the loss
    \item \badgeLsix{L6} Derives the loss from the RLHF objective, explains the implicit reward, discusses $\beta$'s role
    \item \badgeLseven{L7} Discusses DPO variants (IPO, KTO, ORPO), the offline vs.\ online data limitation, and when RLHF is still preferable (e.g., when you need online data collection or reward model reuse across policies)
\end{itemize}

\followups{``What is $\beta$ doing intuitively?'' ``When would you still prefer RLHF over DPO?'' ``What is KTO and how does it differ from DPO?''}
\end{interviewq}

\begin{interviewq}
\textbf{Q: Design the architecture for a customer support ticket classification system with 50 categories, handling multilingual input and evolving category taxonomy.}

\badgeLsix{L6} \quad \badgetype{Architecture Design} \quad \badgefreq{\freqhigh}

\stronganswer I'd approach this as a production classification system with several practical constraints.

\textit{Model selection.} Start with a fine-tuned encoder: mDeBERTa-v3-base for multilingual support (190M params, 100+ languages). Architecture: [CLS] $\to$ Dropout(0.1) $\to$ Linear(768, 50) $\to$ Softmax. An encoder is the right choice because: (1) single forward pass gives sub-10ms latency, (2) fine-tuned encoders beat zero-shot LLMs on domain-specific classification with even 1K labeled examples, (3) cost is $\sim$100$\times$ lower than LLM API calls at scale.

\textit{Handling evolving taxonomy.} Categories change over time (new categories added, old ones merged). Two approaches: (1) \textbf{Retrieval-based}: Encode category descriptions and tickets into the same embedding space (using a dual-encoder like Sentence-BERT), classify by nearest-category. New categories only need a description, no retraining. (2) \textbf{Hierarchical classification}: Group categories into clusters (billing, technical, account). Add a coarse classifier first, then fine-grained. New categories within an existing cluster require less data.

\textit{Data pipeline.} Active learning loop: model classifies incoming tickets with confidence scores. Low-confidence tickets are routed to human agents for labeling. This continuously improves coverage of new patterns and drifting language.

\textit{Evaluation.} Macro-F1 (handles class imbalance), per-category precision/recall, confusion matrix to identify commonly confused pairs. Monitor: classification confidence distribution over time (dropping confidence = distribution drift).

\textit{Fallback strategy.} Tickets below a confidence threshold get routed to an LLM (GPT-4 via API) for zero-shot classification as a safety net. This handles rare categories and novel patterns that the fine-tuned model hasn't learned.

\redflags
\begin{itemize}
    \item Using GPT-4 for all classification---cost prohibitive at scale
    \item Ignoring the evolving taxonomy problem
    \item Not considering class imbalance (some categories are 100$\times$ more common)
\end{itemize}

\interviewtest{Practical system design for NLP, balancing model quality, cost, and maintainability.}

\goodgreat
\begin{itemize}
    \item \badgeLfive{L5} Chooses the right model family (encoder), justifies over LLM
    \item \badgeLsix{L6} Addresses taxonomy evolution, active learning, and evaluation
    \item \badgeLseven{L7} Designs the full production loop (model $\to$ monitoring $\to$ data collection $\to$ retraining), discusses cost analysis (encoder vs.\ LLM API), addresses the cold-start problem for new categories
\end{itemize}

\followups{``How do you handle a new category with zero labeled examples?'' ``What's your retraining cadence?'' ``How do you detect when the model is underperforming?''}
\end{interviewq}

\begin{interviewq}
\textbf{Q: Compare RoPE vs learned positional embeddings vs sinusoidal encodings. Why did RoPE become the dominant choice for modern LLMs?}

\badgeLsix{L6} \quad \badgetype{Conceptual} \quad \badgefreq{\freqmed}

\stronganswer Positional encoding is how transformers understand token order (since self-attention is permutation-invariant without it). The three approaches differ in expressiveness, generalization, and computational cost.

\textit{Sinusoidal (Vaswani et al., 2017):} Fixed, deterministic, uses sine/cosine functions at different frequencies. Advantage: theoretically generalizes to any sequence length. Disadvantage: in practice, it doesn't generalize well beyond training lengths, and the fixed nature limits the model's ability to learn task-specific positional patterns.

\textit{Learned embeddings (BERT, GPT-2):} A trainable embedding lookup table, one vector per position. Advantage: maximum flexibility---the model learns whatever positional patterns are useful. Disadvantage: hard limit at the maximum training length (position 512 for BERT). Cannot generalize to longer sequences at all. The embedding table also adds parameters linear in context length.

\textit{RoPE (Rotary Position Embeddings, Su et al., 2021):} Encodes position by rotating the query and key vectors in 2D subspaces. The rotation angle is proportional to position, so the dot product $q_i^T k_j$ depends only on the relative position $(i - j)$. This gives: (1) \textbf{Relative position awareness}: attention naturally captures relative distances. (2) \textbf{Length extension}: because only relative position matters, trained models can be extended to longer sequences---though zero-shot extrapolation degrades quickly, so in practice extension relies on position interpolation or base-frequency rescaling. (3) \textbf{No extra parameters}: rotations are computed analytically, not learned. (4) \textbf{Decaying long-range dependency}: the paper motivates the design with an upper bound showing inner products decay with relative distance (an implicit locality bias); treat this as the paper's motivation rather than a robust empirical property.

\textit{Why RoPE won:} It combines the generalization of sinusoidal encodings with the relative-position awareness that learned embeddings lack, while adding zero parameters. ALiBi, which replaces position embeddings with a linear distance penalty on attention scores, actually \emph{extrapolates better} to unseen lengths with no tuning. RoPE won anyway on three fronts: better quality at trained lengths, cheap post-hoc context extension (position interpolation, NTK-aware scaling, YaRN, base-frequency rescaling as in Llama 3.1)---which made ALiBi's untuned extrapolation advantage less decisive---and ecosystem momentum once LLaMA standardized it. For modern LLMs with 8K--128K+ contexts, it is RoPE-plus-extension-tricks, not raw RoPE extrapolation, that delivers long context.

\redflags
\begin{itemize}
    \item Not knowing that standard learned embeddings can't handle sequences longer than training
    \item Confusing absolute and relative position encodings
    \item Saying ``sinusoidal encodings generalize perfectly''---they don't in practice
\end{itemize}

\interviewtest{Understanding of positional encoding design choices and their practical implications for long-context models.}

\goodgreat
\begin{itemize}
    \item \badgeLfive{L5} Describes each method and knows RoPE enables length extrapolation
    \item \badgeLsix{L6} Explains the rotation mechanism, relative position property, and why it's parameter-free
    \item \badgeLseven{L7} Discusses NTK-aware scaling, YaRN for extending context beyond training length, and the relationship between RoPE frequency base and effective context window
\end{itemize}

\followups{``What is NTK-aware scaling for RoPE?'' ``How does ALiBi differ from RoPE?'' ``What happens when you test a RoPE model on 4$\times$ the training context length?''}
\end{interviewq}